% ==========================
% APPENDIX A
% ==========================
\chapter{Generating the Sine Function}\label{page:A_appendix}
Appendices include supplementary materials that are not essential for understanding the main content of the thesis, but constitute a valuable extension. These may include, among others, fragments of source code, additional plots, summaries of measurement results, user manuals of the developed system, or auxiliary data used in the analysis.

\section{Source Code}
This appendix presents a sample source code that generates a plot of the \emph{sine} function. The code can be executed in a Python environment (e.g., using the \texttt{NumPy} \cite{NumPyGuide2025} and \texttt{Matplotlib} libraries). The example illustrates how to create simple plots of periodic functions to visualise results in an engineering thesis.

\begin{center}
  \captionsetup{type=listing}
  \inputminted[
    linenos,
    frame=single,
    breaklines,
    breakanywhere
  ]{python}{other/sin.py}
  \caption{Generating a plot of the sine function in Python.}
  \label{lst:sinus_py}
\end{center}

\subsection{Notes on Code Formatting}
Source code should be placed in the \texttt{lstlisting} environment, which preserves formatting and indentation and provides syntax highlighting for the language. If the code is extended, the inserted fragment may be limited, e.g.: \verb|\lstinputlisting[firstline=10,lastline=60]{*.py}|. In engineering theses, it is recommended to include in the appendices only those fragments that are essential for understanding the solution; longer listings are best placed in an online repository, with a reference provided in the text (e.g., \emph{''The source code is available in the \href{https://github.com}{GitHub} repository''}).
